\section{Real-world Application}\label{s:rwa}

% ============================================================
% WORD LIMITS: 2.1 = 250w, 2.2 = 250w
% Excluded: equations, tables, figures, bibliography
% ============================================================

% ============================================================
\subsection{Image Model Deployment}
% ============================================================
% TASK: unpaired image-to-image translation on full frames (not patches).
%       Evaluate in both directions with metrics.
%       ≥10 success + ≥10 failure comparisons. 250 words.
%
% SUGGESTED STRUCTURE:
%
%   [Model] ~50 words
%     "We adopt Contrastive Unpaired Translation (CUT)~\cite{park2020},
%      which trains a generator \(G: A \to B\) using a PatchNCE loss
%      that maximises mutual information between input and output patches,
%      requiring no paired data.  We initialise from the publicly available
%      horse2zebra pretrained weights~\cite{park2020} and fine-tune for
%      5\,+\,3 epochs (constant LR then linear decay) on 500 full frames
%      per domain extracted at the timestamps of 1.1 detections
%      (GPU: NVIDIA Turing, $\sim$X minutes)."
%     Hardware: NVIDIA GeForce RTX 2080 Ti (NCC HPC cluster).
%     NOTE: fill in actual training time once you have it (check slurm logs).
%
%   [Evaluation] ~80 words
%     Metrics: FID (distribution-level quality), KID (unbiased FID),
%     LPIPS (perceptual patch similarity).
%     Report both directions: game→movie (AtoB) and movie→game (BtoA).
%     Example sentence:
%     "Table~\ref{tab:metrics21} reports FID, KID, and LPIPS in both
%      translation directions.  Game-to-movie translation achieves
%      FID\,=\,X / KID\,=\,X, reflecting the larger domain gap from
%      the constrained game palette to the varied cinematic style of
%      movie footage.  Movie-to-game (BtoA) is easier, yielding
%      FID\,=\,X, as the generator maps diverse movie textures to
%      the more homogeneous game rendering."
%
%   [Success cases] ~60 words
%     "Figure~\ref{fig:21_success} shows 10 successful translations.
%      The model consistently transfers film-grain texture and colour
%      grading to game frames, particularly for well-lit frontal shots
%      with clear clothing texture. [Add 1-2 specific visual observations
%      from your actual results, e.g., colour shift, grain, contrast.]"
%
%   [Failure cases] ~60 words
%     "Figure~\ref{fig:21_fail} shows 10 failure cases.
%      Common failures include: (1) background style leaking onto
%      foreground where YOLO bounding boxes were imprecise; (2) colour
%      artefacts on dark game scenes with low texture information;
%      (3) temporal inconsistency visible across consecutive frames
%      as the model has no inter-frame memory."
%
%   [Limitation] ~40 words
%     "The horse2zebra pretrained weights introduce a domain gap:
%      the generator's prior is tuned to animal textures rather than
%      human clothing.  Fine-tuning for only 8 epochs partially corrects
%      this; longer training or a more suitable pretrained base
%      (e.g.\ cityscapes) may improve results."
%
% KEY THINGS TO STATE:
%   - Hardware used and training time (explicitly required for marks)
%   - Both translation directions evaluated
%   - Why these metrics (FID = global distribution, LPIPS = local perceptual)
%   - Honest failure analysis (domain gap from horse2zebra, no temporal coherence)
%
% FIGURES:
%   Fig 4 — Metrics table.
%            Columns: Direction | FID ↓ | KID ↓ | LPIPS ↓
%            Rows: AtoB (game→movie), BtoA (movie→game)
%            Save to: paper/figs/ (as a LaTeX table, not image)
%
%   Fig 5 — Success cases: 2×5 or 5×2 grid, input|output pairs.
%            Pick 10 frames where the style transfer looks convincing.
%            From: output/q2_1/cut_data/testA/ (input)
%                  output/q2_1/<run>/fake_M/ (output)
%            Save to: paper/figs/success_21.pdf
%            Caption: "Successful game-to-movie translations (CUT, 2.1).
%                      Left: input game frame. Right: translated output."
%
%   Fig 6 — Failure cases: same format, 10 bad examples.
%            Save to: paper/figs/fail_21.pdf
%            Caption: "Failure cases: [describe 2-3 specific failure modes]."
%
% LATEX SKELETON:
% \begin{table}[t]
%   \centering
%   \caption{CUT translation metrics (2.1 full-frame model).
%            FID and KID measure distribution similarity; LPIPS measures
%            perceptual patch distance. Lower is better for all.}
%   \label{tab:metrics21}
%   \begin{tabular}{lrrr}
%     \toprule
%     Direction & FID\,↓ & KID\,↓ & LPIPS\,↓ \\
%     \midrule
%     Game $\to$ Movie (AtoB) & -- & -- & -- \\
%     Movie $\to$ Game (BtoA) & -- & -- & -- \\
%     \bottomrule
%   \end{tabular}
% \end{table}


% ============================================================
\subsection{Local (Temporal) Enhancement}
% ============================================================
% TASK: improve 2.1 using local (patch) methods and/or temporal approaches.
%       Submit enhanced video. ≥10 key-frame comparisons. 250 words.
%
% SUGGESTED STRUCTURE:
%
%   [Motivation] ~40 words
%     "The 2.1 model applies style transfer to full frames, mixing
%      foreground human style with background. We address this with
%      patch-level compositing: translating only detected human regions
%      and compositing them onto the original frame."
%
%   [Method overview] ~80 words
%     "For each test frame: (1) YOLO detects human bounding boxes;
%      (2) the 1.2 GCN classifies each patch, discarding `others';
%      (3) kept crops are translated by a CUT model fine-tuned on the
%      1.3-selected patches (game: X, movie: Y patches); (4) translated
%      crops are resized to original bbox dimensions and composited onto
%      the untranslated frame; (5) exponential moving average (EMA)
%      temporal blending (\(\alpha=0.3\)) reduces inter-frame flicker."
%     Fill in actual X, Y counts.
%
%   [Why patch fine-tuning helps] ~40 words
%     "Fine-tuning on human patches (vs.\ full frames) aligns the
%      generator's training distribution with the inference distribution:
%      the model sees the same crop scale and content at both train and
%      test time, reducing identity mapping artefacts."
%
%   [Evaluation — metrics] ~40 words
%     "Table~\ref{tab:metrics22} compares 2.1 and 2.2 metrics on the
%      same test splits.  The enhanced model achieves FID\,=\,X vs.\
%      2.1's Y, a Z\,\% improvement, confirming that patch-level
%      training narrows the domain gap for human regions."
%
%   [Key-frame comparison] ~30 words
%     "Figure~\ref{fig:keyframes} shows 10 key frames: original,
%      2.1 baseline, and 2.2 enhanced.  The patch compositing
%      preserves background fidelity while improving human-region
%      style consistency."
%
%   [Critical evaluation / failure analysis] ~60 words
%     "Remaining failure modes: (1) GCN misclassification of `others'
%      (9.8\% val accuracy) causes occasional style transfer on
%      non-human regions; (2) EMA blending introduces ghosting on fast
%      motion; (3) patch boundary artefacts are visible at bbox edges
%      due to hard compositing without feathering; (4) the horse2zebra
%      pretrained base still limits stylistic range."
%
% KEY THINGS TO STATE:
%   - Hardware: NVIDIA GeForce RTX 2080 Ti (NCC HPC cluster) — state this explicitly
%   - Training time for patch fine-tuning (check slurm logs for wall time)
%   - Explicit comparison to 2.1 (this is the main point)
%   - Both videos submitted (baseline 2.1 and enhanced 2.2)
%   - Honest failure analysis — the marker rewards critical thinking
%
% FIGURES:
%   Fig 7 — Metrics comparison table (2.1 vs 2.2).
%            Columns: Model | FID ↓ | KID ↓ | LPIPS ↓
%            Rows: 2.1 Full-frame CUT, 2.2 Patch CUT + EMA
%
%   Fig 8 — Key-frame comparison (≥10 frames required).
%            Three-column layout: Original | 2.1 Baseline | 2.2 Enhanced
%            Extract representative frames from:
%              output/q2_1/<run>/enhanced_video_21.mp4 (2.1)
%              output/q2_2/<run>/enhanced_model.mp4   (2.2)
%            Save individual frames as: paper/figs/kf_XX_orig.jpg etc.
%            or use a script to make a contact sheet.
%            Caption: "Key frame comparisons: original (left),
%                      2.1 full-frame CUT (centre), 2.2 patch-level
%                      CUT + EMA blending (right)."
%            Save to: paper/figs/keyframes_comparison.pdf
%
%   Fig 9 (optional but strong) — Patch compositing diagram.
%            Annotated figure showing:
%              (a) original frame with YOLO boxes drawn
%              (b) extracted patches
%              (c) translated patches
%              (d) composited result
%            Save to: paper/figs/compositing_pipeline.pdf
%            This is a diagram so word-count exempt and tells the story
%            of 2.2 at a glance.
%
% LATEX SKELETON:
% \begin{figure}[t]
%   \centering
%   \includegraphics[width=\linewidth]{figs/keyframes_comparison.pdf}
%   \caption{Ten key frames from the test video: original (left),
%            2.1 full-frame CUT baseline (centre), and 2.2 patch-level
%            CUT with EMA temporal blending (right).}
%   \label{fig:keyframes}
% \end{figure}
%
% \begin{table}[t]
%   \centering
%   \caption{Metric comparison between 2.1 and 2.2 models.}
%   \label{tab:metrics22}
%   \begin{tabular}{lrrr}
%     \toprule
%     Model & FID\,↓ & KID\,↓ & LPIPS\,↓ \\
%     \midrule
%     2.1 Full-frame CUT   & -- & -- & -- \\
%     2.2 Patch CUT + EMA  & -- & -- & -- \\
%     \bottomrule
%   \end{tabular}
% \end{table}
